\printedpage{24}
% head: A. Grothendieck : Dérivateurs, Ch. II, Édition des Universités de Montpellier II et Paris VII
% foot: Version 8 juillet 2001
\msp{page 37}

\section*{3\quad Return to cofinality}

It is instructive to relate questions of cofinality (for an arrow $u : I' \longrightarrow I$ in Cat) to the derivator $\HOT_{W_0}$, where $W_0 \subseteq \Fl\,\Cat$ is the coarsest of the localizers in Cat satisfying the connectedness axiom Loc 5, i.e.
\[
W_0 = \{ f \in \Fl\,\Cat \mid \pi_0(f) \text{ bijective}\}. \tag*{(3.1)}
\]
One then finds, for $I$ in Cat, a canonical equivalence, 2-functorial in $I$,
\[
\HOT_{W_0}(I) \xrightarrow{\;\sim\;} I^{\wedge} \defeq \uHom(I,(\Ens)), \quad \text{\pl}\ \uHom(I^o,\Ens)\,\text{\texttt{]}} \tag*{(3.2)}
\]
in other words, one has an equivalence
\[
\HOT_{W_0} \simeq \mathbf{D}_{(\Ens)}. \tag*{(3.3)}
\]
When $\HOT_{W_0}$ is defined by localization starting from the category $\Cat(I) \defeq \uHom(I^o,\Cat)$, the arrow (3.2) is the one deduced, by localization, from the arrow
\[
\begin{array}{rcl}
\pi_0^I : \uHom(I^o,\Cat) & \longrightarrow & \uHom(I^o,(\Ens)) \\[1ex]
X & \longmapsto & \pi_0 \circ X,
\end{array} \tag*{(3.4)}
\]
where
\[
\pi_0 : \Cat \longrightarrow (\Ens) \tag*{(3.5)}
\]
is the ``connected components'' functor. If one describes $\HOT_{W_0}(I)$ as a localized category of $\Cat/I$, (3.2) is deduced from the arrow
\[
\left\{
\begin{array}{rcl}
\Cat/I & \longrightarrow & I^{\wedge} = \uHom(I^o,(\Ens)) \\[1ex]
X & \longmapsto & \underline{\pi_0}(X/I) = (i \longmapsto \pi_0(i\backslash X)).
\end{array}
\right. \tag*{(3.6)}
\]

\msp{page 38}
Thus, the familiar functors, for an arrow
\[
u : I' \longrightarrow I
\]
in Cat
\[
\left\{
\; I'^{\wedge} \quad
\begin{array}{c}
\xrightarrow{\;u_!\;} \\[1.5ex]
\xleftarrow{\;u^*\;} \\[1.5ex]
\xrightarrow{\;u_*\;}
\end{array}
\quad I^{\wedge},
\right. \tag*{(3.7)}
\]
\printedpage{25}
% head: A. Grothendieck : Dérivateurs, Ch. II, Édition des Universités de Montpellier II et Paris VII
% foot: Version 8 juillet 2001
are interpreted as the functors of the same name for the derivator $\HOT_{W_0}$, or more directly, for the derivator (which one might mistakenly regard as hoax-like) $\mathbf{D}_{(\Ens)}$.

This gives a meaning to the cohort of cohomological notions associated with a derivator, or more particularly, with a basic localizer $W$ on Cat (giving rise to the derivator $\HOT_W$). We shall designate them by the name of $W_0$-\textit{fourbi} (gear) (instead of $\HOT_{W_0}$-\textit{fourbi}, or $\mathbf{D}_{(\Ens)}$-\textit{fourbi}), in particular for an arrow $u : I' \longrightarrow I$ of Cat such as the notion
\[
\left\{
\begin{array}{lll}
\text{$W_0$-equivalence} & \Bigl(\iff & \text{$\pi_0(u)$ bijective.}\Bigr) \\
\text{$W_0$-aspheric} & \Bigl(\iff & \text{the $I'/i$ for $i \in \Ob I$ are} \\
 & & \text{$0$-connected.}\Bigr) \\
\text{$W_0$-coaspheric} & \Bigl(\iff & \text{the $i\backslash I'$ for $i \in \Ob I$ are} \\
 & & \text{$0$-connected.}\Bigr) \\
\text{$W_0$-proper (\footnotemark[46])} & \Bigl(\iff & \text{the inclusions $I'_i \longrightarrow I'/i$ ($i \in \Ob I$)} \\
 & & \text{are $W_0$-coaspheric, i.e. the fibers} \\
 & & \text{of the colocalized $i'\backslash I' \longrightarrow i\backslash I$} \\
 & & \text{($i' = u(i)$) \pl\ ($i = u(i')$)]} \\
 & & \text{are $0$-connected.}\Bigr) \\
\text{\mspin{page 39}} & & \\
\text{$W_0$-smooth (\footnotemark[47])} & \Bigl(\iff & \text{the inclusions $I'_i \longrightarrow i\backslash I'$} \\
 & & \text{are $W_0$-aspheric, i.e. the fibers} \\
 & & \text{of the localized arrows $I'/i' \longrightarrow I/i$} \\
 & & \text{($i = u(i')$) are $0$-connected.}\Bigr) \\
\text{$W_0$-perfect} & \Bigl(\iff & \text{$W_0$-proper and $W_0$-smooth.}\Bigr) \\
\text{universal $W_0$-equivalence (\footnotemark[48])} & \Bigl(\iff & \text{$\pi_0$ bijective and remains so under any} \\
 & & \text{base change. These are also} \\
 & & \text{the $W_0$-fibrations (or weak $W_0$-fibrations} \\
 & & \text{with $W_0$-aspheric fibers,} \\
 & & \text{i.e. $0$-connected.}\Bigr) \\
\text{$W_0$-fibration} & \Bigl(\iff & \text{$u^* : \Cat/I \longrightarrow \Cat/I'$ transforms} \\
 & & \text{$W_0$-equivalences into $W_0$-equivalences.}\Bigr) \\
\text{weak $W_0$-fibration (\footnotemark[49])} & \Bigl(\iff & \text{the base change} \\
 & & \text{$u^* : \Cat/I \longrightarrow \Cat/I'$ transforms} \\
 & & \text{homotopisms $J_1 \longrightarrow J_2$ into} \\
 & & \text{$W_0$-equivalences - and it suffices to} \\
 & & \text{require it for the inclusion of a} \\
 & & \text{final object (or at one's choice initial) and for} \\
 & & \text{the inclusion of a point in $\boldsymbol{\Delta}^1$.}\Bigr)
\end{array}
\right. \tag*{(3.8)}
\]
\footnotetext[46]{Example: the Cat-cofibrations.}
\footnotetext[47]{Example: the Cat-fibrations.}
\footnotetext[48]{Example: the $W_0$-smooth or $W_0$-proper arrows, with 0-connected fibers.}
\footnotetext[49]{Example: the $W_0$-smooth arrows such that $\pi_0(I'/I) = (i \longmapsto \pi_0(i\backslash I')) = u_!(e_{I'^{\wedge}})$ is locally constant on $I$, i.e. $\underline{\pi_0}(I'/I)$ is an étale \textit{covering} of $I$, or if $u^{\mathrm{opp}}$ satisfies this condition.}

\printedpage{26}
% head: A. Grothendieck : Dérivateurs, Ch. II, Édition des Universités de Montpellier II et Paris VII
% foot: Version 8 juillet 2001
\texttt{[We have a]} \textit{strict} implication (and how)
\[
\text{universal $W_0$-equivalence} \Longrightarrow \text{$W_0$-fibration}
\]
\texttt{[and a]} \textit{strict} implication (and how\,!) (\footnote[50]{But if the fibers are 0-connected, the two notions of $W_0$-fibrations are equivalent and mean that one has a universal $W_0$-equivalence \dots})
\[
\text{$W_0$-fibration} \Longrightarrow \text{weak $W_0$-fibration.}
\]

Let us point out, as an example of the use of these notions, that for a cartesian diagram in Cat
\[
\begin{tikzcd}
I' \arrow[d, "u"'] & J' \arrow[l, "p'"'] \arrow[d, "v"] \\
I & J, \arrow[l, "p"']
\end{tikzcd} \tag*{(3.9)}
\]
the base change arrow for $u_*$, $v_*$ for presheaves is \texttt{[an]} isomorphism if $u$ is $W_0$-proper, or $p$ $W_0$-smooth, and the base change arrow for $u_!$, $v_!$ is \texttt{[an]} isomorphism if $u$ is $W_0$-smooth or $p$ $W_0$ proper (\footnote[51]{Unless I am mistaken, this gives \textit{two characterizations} of $W_0$-proper morphisms, or of $W_0$-smooth morphisms.}). I point out in this connection that these statements, which a priori are formulated for presheaves with values in (Ens), remain valid % [sic] ca -> ça (in footnote 51), missing cedilla
\msp{page 40}
for presheaves with values in any category $\mathcal{M}$ (subject to the existence of the cohomological or homological direct image operations under consideration). This can be seen e.g. by noting that % [sic] operations -> opérations, missing accent
\[
\text{$W_0$-equivalence} \;\Longrightarrow\; \text{$\mathbf{D}_{\mathcal{M}}$-equivalence} \tag*{(3.10)}
\]
for every $\mathcal{U}$-category $\mathcal{M}$, from which it follows that every (\footnote[52]{Not quite correct for \textit{every} ``$W_0$-notion'', one must except the two notions of strong, and weak $W_0$-fibration.}) ``$W_0$-notion'' implies the corresponding $\mathbf{D}_{\mathcal{M}}$-notion, and thus allows one to apply to $\mathbf{D}_{\mathcal{M}}$ the results concerning it. (There is of course a direct, less learned argument, consisting in reducing to the case where $\mathcal{M}$ is small, and then using the fully faithful embedding $\mathcal{M} \hookrightarrow \mathcal{M}^{\wedge}$ resp. $\mathcal{M} \hookrightarrow ((\mathcal{M}^o)^{\wedge})^o$.)

I recall for the record a (third\,?) characterization of the arrows $u : I' \longrightarrow I$ which are $W_0$-proper (resp. $W_0$-smooth): they are those for which the base change functor $u^* : \Cat/I \longrightarrow \Cat/I'$ transforms $W_0$-coaspheric (resp. $W_0$-aspheric) arrows into likewise. This gains in interest, from the fact that the notion of a $W_0$-coaspheric (resp. $W_0$-aspheric) functor is interpreted as that of a cofinal (resp. co-cofinal) functor, cf. below.

\printedpage{27}
% head: A. Grothendieck : Dérivateurs, Ch. II, Édition des Universités de Montpellier II et Paris VII
% foot: Version 8 juillet 2001
\msp{page 41}
The notion of $W_0$-fibration appears extremely strong - much stronger in certain respects than that (say) of $W_\infty$-fibration or of $W_\omega$-fibration. It is equivalent to this (as one sees at once): for every connected component $I'_\alpha$ of $I'$, denoting by $I_\alpha$ the corresponding connected component of $I$, and for every $J_\alpha$ 0-connected and over $I_\alpha$, the fiber product $I'_\alpha \times_{I_\alpha} J_\alpha$ is again 0-connected. (That looks like those stories of total 0-connectedness of a topos\,!) But this also means that $I'_\alpha \longrightarrow I_\alpha$ is a universal $W_0$-equivalence. This implies that its fibers are 0-connected, and the converse is true if $u : I' \longrightarrow I$ is $W_0$-proper or $W_0$-smooth. Thus, the essential condition for a $W_0$-fibration is that the fibers of the induced arrows $I'_\alpha \longrightarrow I_{u(\alpha)}$ \pl\ $I'_\alpha \longrightarrow I_\alpha$\texttt{]} be 0-connected. In fact, one has this

\noindent\textbf{Proposition 3.1.} \textit{Let $u : I' \longrightarrow I$ be an arrow in} Cat.
\begin{itemize}
\item[\textit{a)}] \textit{For it to be a universal $W_0$-equivalence, it is necessary and sufficient that for every $\boldsymbol{\Delta}^1 \longrightarrow I$, i.e. every arrow $\alpha$ of $I$, the category $I' \times_I \boldsymbol{\Delta}^1 \defeq I'|\alpha$ be $0$-connected, or again that the fibers of $I$ be $0$-connected, and that for every} nonidentity \textit{arrow $\alpha$ of $I$, $I'|\alpha$ be $0$-connected.}
\item[\textit{b)}] \textit{For $u : I' \longrightarrow I$ to be a $W_0$-fibration, it is necessary and sufficient that for every connected component $I'_{\alpha'}$ of}
\msp{page 42}
\textit{$I'$, putting $\alpha = u(\alpha')$, the induced arrow $I'_{\alpha'} \longrightarrow I_\alpha$ be a universal $W_0$-equivalence (cf. criterion a) for this notion).}
\item[\textit{c)}] \textit{For it to be a weak $W_0$-fibration, it is} necessary \textit{that $\forall\, i \in I$, the inclusion $I'_i \hookrightarrow i\backslash I'$ be a $W_0$-equivalence, i.e. $\pi_0(I'_i) \xrightarrow{\;\sim\;} \pi_0(i\backslash I')$ be bijective,} and \textit{that $i \longmapsto \pi_0(i\backslash I')$ be a locally constant presheaf on $I$ (i.e. $\underline{\pi_0}(I'/I)$ be an étale} covering \textit{of $I$). This condition is also sufficient, if $u$ is $W_0$-smooth. Moreover (for the record) $u$ is a weak $W_0$-fibration if and only if $u^{\mathrm{opp}}$ is.}
\end{itemize}

The greatest difference between the notion of $W_0$-fibration, and that of weak $W_0$-fibration, seems to me that $W_0$-perfect morphisms (i.e. $W_0$-smooth and $W_0$-proper), and a fortiori the $W$-perfect ones for a basic localizer $W \subseteq W_0$, are necessarily weak $W$-fibrations, but (\textit{even if $W = W_\omega$,} the finest of all basic localizers) are not $W$-fibrations. The most striking example in this respect is probably that of étale coverings, which are $W_\omega$-perfect (hence
\msp{page 43}
weak $W_\omega$-fibrations and a fortiori weak $W_0$-fibrations), but not $W_0$-fibrations\,; for in the case where $I'$, $I$ are 0-connected, the fibers would have to be 0-connected (cf. b) above), i.e. $u$ would have to be an isomorphism\,!

A nice property of $W_0$-smooth (resp. $W_0$-proper) morphisms, is that they are \textit{open} (resp. closed) maps for the underlying topological spaces, which indeed means that the local fibers, i.e. of the arrows localized above $I'/i' \longrightarrow I/i$ (resp. colocal, i.e. the arrows colocalized below $i'\backslash I' \longrightarrow i\backslash I$), are surjective, i.e. with nonempty fibers (which is implied by the condition that these be 0-connected).

\printedpage{28}
% head: A. Grothendieck : Dérivateurs, Ch. II, Édition des Universités de Montpellier II et Paris VII
% foot: Version 8 juillet 2001
In all this I have still hardly spoken of $W_0$-aspheric resp. coaspheric arrows, which nevertheless are here a main motivation for introducing $W_0$\,!

\noindent\textbf{Proposition 3.2.} \textit{Equivalent conditions on the arrow $u : I' \longrightarrow I$ in} Cat:
\begin{itemize}
\item[a)] \textit{For every $\mathcal{U}$-category $\mathcal{M}$ and every functor $F : I \longrightarrow \mathcal{M}$, the canonical arrow (provided that small direct limits are representable in $\mathcal{M}$),}
\[
\varinjlim_{I'} u^*(F) \longrightarrow \varinjlim_{I} F \tag*{(3.2.1)}
\]
\msp{page 44}
\textit{is} \texttt{[an]} \textit{isomorphism.}
\item[a$'$)] \textit{For every $\mathcal{U}$-category $\mathcal{M}$ in which small inverse limits are representable and every $F \in \mathcal{M}(I) = \uHom(I^o,\mathcal{M})$, i.e. every presheaf on $I$ with values in $\mathcal{M}$, the canonical arrow}
\[
\varprojlim_{I} F \longrightarrow \varprojlim_{I'} u^*(F) \tag*{(3.2.2)}
\]
\textit{is} \texttt{[an]} \textit{isomorphism.}
\item[a$''$)] (\footnote[53]{This is a condition on morphisms of topoi, namely that if $F$ \texttt{[is]} in $\mathcal{X}^{\wedge}$, one has
\[
H^0(\mathcal{X},F) \xrightarrow{\;\sim\;} H^0(\mathcal{X}',u^*F).
\]
There is \texttt{[a]} similar notion for the $H^i$ ($i \leq n$), to be brought out. Notion of $W_0$-coaspheric morphisms of topoi\,! It is in the nature of a relative 0-connectedness condition, dual to the $W_0$-aspheric notion, which simply means that for $U$ 0-connected in $U^{\wedge}$ \pl\ $\mathcal{X}^{\wedge}$\texttt{]}, $f^*(U)$ \pl\ $u^*(U)$\texttt{]} is 0-connected. Likewise, one probably has notions (``dual'' to the preceding ones) of $W_n$-coasphericity (well understood in principle if $\mathcal{X}$ is locally $n$-aspheric) (The $W_n$-coasphericity, being that mentioned above.)}) \textit{As in a$'$, with $\mathcal{M} = (\Ens)$, i.e. $F \in \Ob I^{\wedge}$.}~(\footnote[54]{\textbf{NB} One can introduce
\begin{itemize}
\item[a$'''$)] As in a), but with $\mathcal{M} = (\Ens)$.
\item[a$^{IV}$)] As in a), with the canonical functor $I' \longrightarrow I'^{\wedge}$.
\end{itemize}}).
\item[b)] \textit{The arrow $u$ is $W_0$-coaspheric, i.e. for every $i \in I$, $i\backslash I'$ is $0$-connected.}
\end{itemize}
The equivalence of a, a$'$, a$''$ is $\pm$ tautological in terms of the definitions (cf. SGA 4, I, 8\dots). The form a$''$ reads, putting $\mathbf{D} = \HOT_{W_0}$, that $\forall\, F \in \mathbf{D}(I)$, one wants
\[
\underbrace{H^{\bullet}_{\mathbf{D}}(I,F)}_{\defeq p_{I,*}(F)} \longrightarrow \underbrace{H^{\bullet}_{\mathbf{D}}(I',u^*F)}_{\defeq p_{I',*}(F)} \qquad \text{\pl}\quad \underbrace{H^{\bullet}_{\mathbf{D}}(I',u^*F)}_{\defeq p_{I',*}(u^*F)}\ \text{\texttt{]}} \tag*{(3.2.3)}
\]
($p_I : I \longrightarrow e$ and $p_{I'} : I' \longrightarrow e$ the canonical functors) to be \texttt{[an]} isomorphism. One knows that this is equivalent to $u$ $\mathbf{D}$-coaspheric. The proof is an amusing game between $u_!$, $u^*$, $u_*$ all at once. One must express that the Hom from an object $M$ of $\mathbf{D}(e)$ ($=$ (Ens) in this case) into the arrow (3.2.3) is bijective. Here, it % [sic] occurence -> occurrence, French spelling